\chapter*{Revision Map and Capstone Practice}
\addcontentsline{toc}{chapter}{Revision Map and Capstone Practice}
\markboth{Revision Map and Capstone Practice}{}

\section*{The five exam-solving habits}
Most Operating Systems questions become manageable if you use a repeatable method:
\begin{enumerate}
  \item \textbf{Name the state.} Write the queue, matrices, frame contents, free blocks, or directory nodes before applying a policy.
  \item \textbf{Name the event.} Arrival, dispatch, quantum expiry, resource request, memory request, page reference, or file-block allocation.
  \item \textbf{Name the policy.} FCFS, SJF, Priority, RR, Banker's safety rule, first/best/worst fit, FIFO/LRU/LFU.
  \item \textbf{Update exactly once.} Change the state according to the event and policy; do not mentally skip intermediate states.
  \item \textbf{Compute metrics last.} Waiting time, turnaround, page faults, fragmentation, or safe sequence should follow from the completed trace.
\end{enumerate}

\begin{mltable}
\begin{tabularx}{\linewidth}{@{}>{\bfseries\leavevmode\color{MLNavy}}p{0.24\linewidth}p{0.30\linewidth}X@{}}
\toprule
\rowcolor{MLPaleBlue!92}
Question type & State to draw & Final evidence \\
\midrule
CPU scheduling & Ready queue + Gantt chart & Completion, turnaround, waiting, response \\
Banker's algorithm & Allocation, Max, Need, Work, Finish & Safe sequence or reason none exists \\
Memory fit & Free-block table & Chosen block + fragment for each request \\
Page replacement & Frame table after every reference & Hit/fault count and victim sequence \\
File allocation & Free/used block map & Valid block mapping with conflicts rejected \\
\bottomrule
\end{tabularx}
\end{mltable}

\section*{Capstone practical}
Build one command-line simulator that lets a user choose between CPU scheduling, memory allocation, or page replacement. Separate \emph{policy} from \emph{state}: for example, one common frame-state representation can be reused with FIFO, LRU, and LFU victim-selection functions.

\begin{miniproject}
\textbf{Minimum features.}
\begin{itemize}
  \item Scheduling: FCFS, SJF, Priority, RR with a printed Gantt trace and average waiting/turnaround time.
  \item Memory: first, best, worst fit with fragmentation results.
  \item Virtual memory: FIFO, LRU, LFU with per-reference frames, hits, and faults.
  \item Input validation and deterministic tie-breaking rules.
  \item A short README explaining assumptions and how to compile/run the program.
\end{itemize}
\end{miniproject}

\begin{decisionmemo}
When two algorithms give different answers, do not ask only ``which is better?'' Ask which objective matters: responsiveness, throughput, fairness, memory utilization, implementation overhead, locality, or predictable behavior.
\end{decisionmemo}

\section*{Final checklist}
\begin{quickcheck}
\begin{enumerate}
  \item Can you explain the parent/child return values of \texttt{fork()} without running code?
  \item Can you derive waiting and turnaround time from a Gantt chart?
  \item Can you distinguish a race condition from deadlock?
  \item Can you compute Need and search for a Banker's safe sequence?
  \item Can you distinguish internal from external fragmentation?
  \item Can you trace FIFO, LRU, and LFU one reference at a time?
  \item Can you explain why directory hierarchy is not the same thing as physical disk layout?
  \item Can you compare text/binary records and contiguous/linked/indexed allocation?
\end{enumerate}
\end{quickcheck}
